
\documentclass{article}
\usepackage{colortbl}
\usepackage{makecell}
\usepackage{multirow}
\usepackage{supertabular}

\begin{document}

\newcounter{utterance}

\centering \large Interaction Transcript for game `wordle', experiment `high\_frequency\_words\_no\_clue\_no\_critic', episode 4 with glm{-}5.2.
\vspace{24pt}

{ \footnotesize  \setcounter{utterance}{1}
\setlength{\tabcolsep}{0pt}
\begin{supertabular}{c@{$\;$}|p{.15\linewidth}@{}p{.15\linewidth}p{.15\linewidth}p{.15\linewidth}p{.15\linewidth}p{.15\linewidth}}
    \# & \multicolumn{2}{c}{Player} && \multicolumn{2}{c}{Game Master} \\
    \hline

    \theutterance \stepcounter{utterance}  
    & & & \multicolumn{4}{p{0.6\linewidth}}{
        \cellcolor[rgb]{0.9,0.9,0.9}{
            \makecell[{{p{\linewidth}}}]{
                \texttt{\tiny{[P1$\langle$GM]}}
                \texttt{Du är en språkmagiker som gillar att gissa ord genom att använda de givna reglerna.} \\
\\ 
\texttt{Välkommen till Wordle! Du har sex försök på dig att gissa målordet, ett giltigt svenskt ord med fem gemena bokstäver. Använd taggarna "förklaring:" och "gissning:" för att ge en kortfattad förklaring till varje gissning.} \\
\\ 
\texttt{Till exempel, om din gissning är "fråga", ska ditt svar vara} \\
\texttt{förklaring: detta är ett vanligt svenskt ord med fem bokstäver, och jag börjar min gissning med detta ord.} \\
\texttt{gissning: fråga} \\
\\ 
\texttt{Efter varje gissning kommer ditt svar att valideras, och du kommer att få återkoppling som visar vilka bokstäver som är korrekta (gröna), vilka bokstäver som är korrekta men på fel position (gula), och vilka bokstäver som är felaktiga (röda). Denna återkoppling kan vara användbar för att avgöra vilka bokstäver du ska ta med eller utesluta i din nästa gissning.} \\
\\ 
\texttt{Till exempel kan återkopplingen för "fråga" vara:} \\
\texttt{gissning\_återkoppling: f<gul> r<gul> å<grön> g<gul> a<röd>} \\
\\ 
\texttt{Förklaringen ska innehålla detaljer om hur gissning\_återkoppling används för att komma fram till en ny gissning.} \\
\\ 
\texttt{Generera ingen annan text.} \\
\\ 
\texttt{Låt oss börja med din första gissning.} \\
            }
        }
    }
     \\ \\

    \theutterance \stepcounter{utterance}  
    & \multicolumn{4}{p{0.6\linewidth}}{
        \cellcolor[rgb]{0.561,0.737,0.561}{
            \makecell[{{p{\linewidth}}}]{
                \texttt{\tiny{[P1$\rangle$GM]}}
                \texttt{förklaring: detta är ett vanligt svenskt ord med fem bokstäver, och jag börjar min gissning med detta ord.} \\
\texttt{gissning: fråga} \\
            }
        }
    }
    & & \\ \\

    \theutterance \stepcounter{utterance}  
    & & & \multicolumn{2}{p{0.3\linewidth}}{
        \cellcolor[rgb]{0.9,0.9,0.9}{
            \makecell[{{p{\linewidth}}}]{
                \texttt{\tiny{[GM$|$GM]}}
                \texttt{attempts = 1} \\
\texttt{target\_word = dator} \\
\texttt{guess = fråga} \\
\texttt{guess\_feedback = f<red> r<yellow> å<red> g<red> a<yellow>} \\
            }
        }
    }
    & & \\ \\

    \theutterance \stepcounter{utterance}  
    & & & \multicolumn{4}{p{0.6\linewidth}}{
        \cellcolor[rgb]{0.9,0.9,0.9}{
            \makecell[{{p{\linewidth}}}]{
                \texttt{\tiny{[P1$\langle$GM]}}
                \texttt{gissning\_återkoppling: f<röd> r<gul> å<röd> g<röd> a<gul>} \\
\\ 
\texttt{Ge ditt svar endast i detta format.} \\
\texttt{förklaring: detaljer} \\
\texttt{gissning: ord} \\
            }
        }
    }
     \\ \\

    \theutterance \stepcounter{utterance}  
    & \multicolumn{4}{p{0.6\linewidth}}{
        \cellcolor[rgb]{0.561,0.737,0.561}{
            \makecell[{{p{\linewidth}}}]{
                \texttt{\tiny{[P1$\rangle$GM]}}
                \texttt{förklaring: 'f', 'å' och 'g' är felaktiga bokstäver och utesluts. 'r' finns i ordet men inte på position 2, och 'a' finns i ordet men inte på position 5. Jag provar ett nytt ord med 'r' och 'a' på andra positioner.} \\
\texttt{gissning: pilar} \\
            }
        }
    }
    & & \\ \\

    \theutterance \stepcounter{utterance}  
    & & & \multicolumn{2}{p{0.3\linewidth}}{
        \cellcolor[rgb]{0.9,0.9,0.9}{
            \makecell[{{p{\linewidth}}}]{
                \texttt{\tiny{[GM$|$GM]}}
                \texttt{attempts = 2} \\
\texttt{target\_word = dator} \\
\texttt{guess = pilar} \\
\texttt{guess\_feedback = p<red> i<red> l<red> a<yellow> r<green>} \\
            }
        }
    }
    & & \\ \\

    \theutterance \stepcounter{utterance}  
    & & & \multicolumn{4}{p{0.6\linewidth}}{
        \cellcolor[rgb]{0.9,0.9,0.9}{
            \makecell[{{p{\linewidth}}}]{
                \texttt{\tiny{[P1$\langle$GM]}}
                \texttt{gissning\_återkoppling: p<röd> i<röd> l<röd> a<gul> r<grön>} \\
\\ 
\texttt{Ge ditt svar endast i detta format.} \\
\texttt{förklaring: detaljer} \\
\texttt{gissning: ord} \\
            }
        }
    }
     \\ \\

    \theutterance \stepcounter{utterance}  
    & \multicolumn{4}{p{0.6\linewidth}}{
        \cellcolor[rgb]{0.561,0.737,0.561}{
            \makecell[{{p{\linewidth}}}]{
                \texttt{\tiny{[P1$\rangle$GM]}}
                \texttt{förklaring: 'p', 'i' och 'l' är felaktiga. 'r' är korrekt på position 5. 'a' är korrekt men på fel position (inte position 4). Jag letar nu efter ett ord som slutar på 'r' och har 'a' på en annan position.} \\
\texttt{gissning: svart} \\
            }
        }
    }
    & & \\ \\

    \theutterance \stepcounter{utterance}  
    & & & \multicolumn{2}{p{0.3\linewidth}}{
        \cellcolor[rgb]{0.9,0.9,0.9}{
            \makecell[{{p{\linewidth}}}]{
                \texttt{\tiny{[GM$|$GM]}}
                \texttt{attempts = 3} \\
\texttt{target\_word = dator} \\
\texttt{guess = svart} \\
\texttt{guess\_feedback = s<red> v<red> a<yellow> r<yellow> t<yellow>} \\
            }
        }
    }
    & & \\ \\

    \theutterance \stepcounter{utterance}  
    & & & \multicolumn{4}{p{0.6\linewidth}}{
        \cellcolor[rgb]{0.9,0.9,0.9}{
            \makecell[{{p{\linewidth}}}]{
                \texttt{\tiny{[P1$\langle$GM]}}
                \texttt{gissning\_återkoppling: s<röd> v<röd> a<gul> r<gul> t<gul>} \\
\\ 
\texttt{Ge ditt svar endast i detta format.} \\
\texttt{förklaring: detaljer} \\
\texttt{gissning: ord} \\
            }
        }
    }
     \\ \\

    \theutterance \stepcounter{utterance}  
    & \multicolumn{4}{p{0.6\linewidth}}{
        \cellcolor[rgb]{0.561,0.737,0.561}{
            \makecell[{{p{\linewidth}}}]{
                \texttt{\tiny{[P1$\rangle$GM]}}
                \texttt{förklaring: 's' och 'v' är felaktiga. 'r' är inte på position 4. 'a' är inte på position 3. 't' är inte på position 5. Vi vet att ordet slutar på 'r' från föregående gissning, så 't' måste vara på en annan plats. Jag provar ett ord som kombinerar dessa ledtrådar.} \\
\texttt{gissning: matar} \\
            }
        }
    }
    & & \\ \\

    \theutterance \stepcounter{utterance}  
    & & & \multicolumn{2}{p{0.3\linewidth}}{
        \cellcolor[rgb]{0.9,0.9,0.9}{
            \makecell[{{p{\linewidth}}}]{
                \texttt{\tiny{[GM$|$GM]}}
                \texttt{attempts = 4} \\
\texttt{target\_word = dator} \\
\texttt{guess = matar} \\
\texttt{guess\_feedback = m<red> a<green> t<green> a<red> r<green>} \\
            }
        }
    }
    & & \\ \\

    \theutterance \stepcounter{utterance}  
    & & & \multicolumn{4}{p{0.6\linewidth}}{
        \cellcolor[rgb]{0.9,0.9,0.9}{
            \makecell[{{p{\linewidth}}}]{
                \texttt{\tiny{[P1$\langle$GM]}}
                \texttt{gissning\_återkoppling: m<röd> a<grön> t<grön> a<röd> r<grön>} \\
\\ 
\texttt{Ge ditt svar endast i detta format.} \\
\texttt{förklaring: detaljer} \\
\texttt{gissning: ord} \\
            }
        }
    }
     \\ \\

    \theutterance \stepcounter{utterance}  
    & \multicolumn{4}{p{0.6\linewidth}}{
        \cellcolor[rgb]{0.561,0.737,0.561}{
            \makecell[{{p{\linewidth}}}]{
                \texttt{\tiny{[P1$\rangle$GM]}}
                \texttt{förklaring: 'm' och den andra 'a' är felaktiga. 'a' är korrekt på position 2, 't' är korrekt på position 3 och 'r' är korrekt på position 5. Det saknas en bokstav på position 1 och 4, men position 4 är inte 'a'. Jag provar med 'e' på position 4.} \\
\texttt{gissning: later} \\
            }
        }
    }
    & & \\ \\

    \theutterance \stepcounter{utterance}  
    & & & \multicolumn{2}{p{0.3\linewidth}}{
        \cellcolor[rgb]{0.9,0.9,0.9}{
            \makecell[{{p{\linewidth}}}]{
                \texttt{\tiny{[GM$|$GM]}}
                \texttt{attempts = 5} \\
\texttt{target\_word = dator} \\
\texttt{guess = later} \\
\texttt{guess\_feedback = l<red> a<green> t<green> e<red> r<green>} \\
            }
        }
    }
    & & \\ \\

    \theutterance \stepcounter{utterance}  
    & & & \multicolumn{4}{p{0.6\linewidth}}{
        \cellcolor[rgb]{0.9,0.9,0.9}{
            \makecell[{{p{\linewidth}}}]{
                \texttt{\tiny{[P1$\langle$GM]}}
                \texttt{gissning\_återkoppling: l<röd> a<grön> t<grön> e<röd> r<grön>} \\
\\ 
\texttt{Ge ditt svar endast i detta format.} \\
\texttt{förklaring: detaljer} \\
\texttt{gissning: ord} \\
            }
        }
    }
     \\ \\

    \theutterance \stepcounter{utterance}  
    & \multicolumn{4}{p{0.6\linewidth}}{
        \cellcolor[rgb]{0.561,0.737,0.561}{
            \makecell[{{p{\linewidth}}}]{
                \texttt{\tiny{[P1$\rangle$GM]}}
                \texttt{förklaring: 'l' och 'e' är felaktiga. Vi vet att ordet har mönstret \_ a t \_ r. Eftersom 'a' inte får vara på position 4 (från tidigare feedback), och vi har uteslutit många bokstäver, kan 'a' vara på position 1.} \\
\texttt{gissning: katar} \\
            }
        }
    }
    & & \\ \\

    \theutterance \stepcounter{utterance}  
    & & & \multicolumn{2}{p{0.3\linewidth}}{
        \cellcolor[rgb]{0.9,0.9,0.9}{
            \makecell[{{p{\linewidth}}}]{
                \texttt{\tiny{[GM$|$GM]}}
                \texttt{Error: The guessed word is not a valid word for this game.} \\
            }
        }
    }
    & & \\ \\

    \theutterance \stepcounter{utterance}  
    & & & \multicolumn{4}{p{0.6\linewidth}}{
        \cellcolor[rgb]{0.9,0.9,0.9}{
            \makecell[{{p{\linewidth}}}]{
                \texttt{\tiny{[P1$\langle$GM]}}
                \texttt{Din gissning är inte ett giltigt ord för det här spelet. Vänligen försök igen.} \\
\\ 
\texttt{Ge ditt svar endast i detta format.} \\
\texttt{förklaring: detaljer} \\
\texttt{gissning: ord} \\
            }
        }
    }
     \\ \\

    \theutterance \stepcounter{utterance}  
    & \multicolumn{4}{p{0.6\linewidth}}{
        \cellcolor[rgb]{0.561,0.737,0.561}{
            \makecell[{{p{\linewidth}}}]{
                \texttt{\tiny{[P1$\rangle$GM]}}
                \texttt{förklaring: 'l' och 'e' var felaktiga. Mönstret är \_ a t \_ r. Eftersom 'a' är fel på position 4, testar jag med 's' på position 4.} \\
\texttt{gissning: hatar} \\
            }
        }
    }
    & & \\ \\

    \theutterance \stepcounter{utterance}  
    & & & \multicolumn{2}{p{0.3\linewidth}}{
        \cellcolor[rgb]{0.9,0.9,0.9}{
            \makecell[{{p{\linewidth}}}]{
                \texttt{\tiny{[GM$|$GM]}}
                \texttt{attempts = 6} \\
\texttt{target\_word = dator} \\
\texttt{guess = hatar} \\
\texttt{guess\_feedback = h<red> a<green> t<green> a<red> r<green>} \\
            }
        }
    }
    & & \\ \\

    \theutterance \stepcounter{utterance}  
    & & & \multicolumn{2}{p{0.3\linewidth}}{
        \cellcolor[rgb]{0.9,0.9,0.9}{
            \makecell[{{p{\linewidth}}}]{
                \texttt{\tiny{[GM$|$GM]}}
                \texttt{game\_result = LOSS} \\
            }
        }
    }
    & & \\ \\

\end{supertabular}
}

\end{document}
